% Generated from Vault/LaTeX/3-fixture.md (version 1.0, last updated ?). Edit the note: this file is rewritten from it on every save.
% !TEX root = main.tex
\startwordcount{text}
\section{Text and Emphasis}
\label{sec:text}
\printwordcount{text}

% P1.1 | a paragraph label comment, not counted 
A plain paragraph of running text with a one-to-one match, a state-of-the-art estimator, numbers such as 27, 3.5 and 1,000, a share of 50\% and 7\,\%, a (bracketed aside) and a sentence that ends here. Words separated only by a tie: Mr.~Smith and 372\,ms, and a thin space in 10\,000 people.

Emphasis takes \emph{emphasised words}, \textit{italic words}, \textbf{bold words} and \textbf{*bold italic words}*, \texttt{code\_name--with-dashes} and \texttt{two words}, \textcolor{red}{red words here}, and a \href{https://example.org/a-b_c}{link with text} to a page.

Quotes: ``curly double quotes'', `curly single quotes', it's, the users' choice, "straight double" and 'straight single' marks.

Dashes: pages 4--5, an en dash – here, an em dash — there, joined---up, spaced -- out, a lone - hyphen, and well-known self-evident co-operation.

Characters: café, naïve, Łukasz, Ångström, a×b, 5 × 6, x ≈ y, and an ellipsis… then on.

Very long words that may be hyphenated at the ends of lines: internationalisation, characterisation, incomprehensibilities, counterrevolutionaries, institutionalisation, electroencephalographically, uncharacteristically, telecommunications, interdisciplinarity, compartmentalisation, misunderstandings, disproportionately, representativeness, overgeneralisation, incontrovertibility.

% A comment line with several words that are never counted 
A line with an end comment %c not counted
and its continuation.

\subsection{Citations}
\label{sec:text:cites}

A citation \cite{one}. One with a page \cite[p.~27]{one}. A range \cite[pp.~4–5]{two}. Two keys \cite{one, two}. Three keys \cite{one, two, three}. A prefix \cite[see][p.~2]{two}. A locator on the first key \cites[p.~3]{one}{two}. A prefix in the middle \cites{one}[see][p.~4]{two}{three}. Attached to a word\cite{one} and before a stop \cite{three}.

As \citet{one} argues, and as \citet[p.~4]{two} shows, the claim holds. \Citet{three} agree. \Citet{corp} warns. \Citet{prefix} replies. \Citet{edited} collects. \Citet{accent} compare. \Citet{particle} notes. \Citet{others} found it.

Raw forms: \citet{one,two} and \citep{three} and \citeauthor{two} and \textcite{three} and \parencite[p.~9]{one} and \Citet{accent}.

\subsection{Glossary and References}
\label{sec:text:refs}

Glossary terms: \gls{FNRATE}, \glspl{FNRATE}, \gls{knowngood}, \Gls{knowngood}, \glspl{knowngood}, \gls{W1}, \glspl{W1}, \gls{shift}, and \glslink{FNRATE}{the signal error} once.

References: Section~\ref{sec:text:cites}, Sections~\ref{sec:text:cites} and \ref{sec:text:lists}, Table~\ref{tab:demo}, Figure~\ref{fig:demo}, Equation~\ref{eq:demo}, Listing~\ref{lst:bench}, \rqref{1} and \hypref{2a}, and (Section~\ref{sec:text:cites}) in brackets.

\subsection{Lists and Footnotes}
\label{sec:text:lists}

A sentence with a footnote\footnote{A footnote of five words.} and another one\footnote{A second footnote, citing \citet{one} with \emph{emphasis}.} in it.

\begin{itemize}
\item A bullet item of six words
\item A second bullet item
	\item A nested bullet item
		\item A third-level bullet
			\item A fourth-level bullet
\end{itemize}

\begin{enumerate}
\item A numbered item
\item A second numbered item
	\item A nested numbered item
		\item A third-level numbered item
			\item A fourth-level numbered item
\end{enumerate}

The labelled list:

\begin{enumerate}
  \item[\textbf{RQ1}] \hypertarget{rq1}{}Does the first question hold?
  \item[\textbf{H2a}] \hypertarget{hyp2a}{}The second hypothesis, in \emph{italics}.
\end{enumerate}

\phantomsection\label{sec:text:anchor}
A paragraph ending in a label.

\paragraph{A Run-in Heading} Its text follows on the same line.

\startwordcount{floats}
\section{Maths, Floats and Code}
\label{sec:floats}
\printwordcount{floats}

Inline maths $x^2 + y$ in a sentence, an $n$-fold rise, a 10$\times$ speed-up, $\alpha$-level tests, a value of $5$\% and $p<0.05$.

\[
a = b + c
\]

The display above has no number.

\begin{equation}
E = m c^2
  \label{eq:demo}
\end{equation}
The equation above is numbered.

\begin{align}
  a &= b + c \\
  d &= e \nonumber \\
  f &= g \tag{T}
\end{align}

\begin{align*}
  a &= b \\
  c &= d
\end{align*}

\begin{gather}
  a = b \\
  c = d
\end{gather}

\begin{multline}
  a + b + c \\
  = d + e
\end{multline}

\begin{equation*}
  x = y
\end{equation*}

A table follows.

\begin{table}[htbp]
\centering
\caption{A table caption of several words, citing \cite{one}.}
\label{tab:demo}
\small
\begin{tabular}{@{}lc@{}}
\toprule
Column one & Column two \\
\midrule
table words & more table words \\
a \gls{FNRATE} cell & \cite{one} \\
\bottomrule
\end{tabular}
\end{table}
\begin{figure}[htbp]
\centering
\resizebox{\ifdim\width>\linewidth\linewidth\else\width\fi}{!}{%
\begin{tikzpicture}
\node (a) {picture words};
\end{tikzpicture}
}
\caption{A figure caption of several words.}
\label{fig:demo}
\end{figure}
\begin{handtyped}[htbp]
\handtypedcode{bench}
\caption{A listing caption of several words}
\label{lst:bench}
\end{handtyped}
\begin{verbatim}
verbatim code block words
\end{verbatim}

\FloatBarrier

{\wordcountoff These words are left out of the count.} Counted again after the group.

\subsection*{Unnumbered Heading}
\phantomsection\addcontentsline{toc}{subsection}{Unnumbered Heading}
\label{sec:floats:unnumbered}

Words under an unnumbered heading.

% =====================================================================

A paragraph after a divider.
\stopwordcount
